\chapter{Conclusion Générale}

% Rappelle de la problématique 

Tout au long de ce travail, notre réflexion s'est structurée autour d'une interrogation centrale : dans quelle mesure les choix inhérents à la \og mise en données \fg{} des sources conditionnent-ils l'efficacité heuristique des modèles d'apprentissage profond ? Autrement dit, comment les décisions prises en amont par l'expérimentateur déterminent-elles la capacité de la machine à produire des connaissances, ou des fractions de connaissances, qui soient véritablement pertinentes pour l'historien ?

Pour répondre à cette question, il nous a fallu parcourir l'ensemble de la chaîne de production algorithmique, en refusant de considérer le réseau de neurones comme une \og boîte noire \fg{} déconnectée de la matérialité du travail érudit. Nous sommes d'abord partis d'une mise en perspective historique et technique, afin de comprendre comment s'est noué, autour de la numérisation des collections, cette collaboration entre sciences informatiques et sciences historiques. Nous avons ensuite plongé au cœur de la fabrique des annotations, démontrant que la constitution d'un corpus, la modélisation d'une ontologie (notamment via les principes \gls{agc}) et la rigueur d'une campagne d'annotation constituent des actes scientifiques à part entière. Enfin, l'entraînement et l'évaluation technique des modèles nous ont permis de confronter l'efficacité algorithmique brute à la validité historique des résultats, révélant ainsi les impacts directs de nos choix de modélisation et les conseils pour améliorer nos résultats. 

% Bilan de la réflexion

Dans une première partie, nous avons retracé l'évolution du traitement automatique des documents historiques, des premiers réseaux de neurones appliqués au secteur postal jusqu'aux modèles multimodaux actuels. Ensuite, nous avons démontré que l'application de la vision artificielle aux documents historiques relève moins de la performance algorithmique que d'un enjeu méthodologique centré sur la structuration des données. Au-delà des tâches techniques (classification, détection, segmentation) et des métriques d'évaluation informatiques (Précision, Rappel, Score F1), la scientificité de la démarche historique repose sur la construction rigoureuse de la « vérité de terrain » et de son ontologie d'annotation. Cette évolution technique transforme profondément la pratique de l'historien, le poussant à délaisser le travail isolé au profit d'une recherche collective dépendante d'infrastructures partagées. Dans ce contexte, l'adoption de standards ouverts et interopérables comme le protocole \gls{iiif}, au cœur d'outils fédérateurs comme AIKON, devient indispensable pour manipuler les sources à distance, garantir la reproductibilité des résultats et favoriser l'émergence d'une véritable science ouverte en histoire. Enfin, nous avons présenté AIKON et l'écosystème de la recherche et développement en vision appliquée aux documents historiques ainsi que l'apparition des nouvelles méthodes que les historiens peuvent appliquer. En rendant la vision par ordinateur accessible sans prérequis technique, la plateforme permet aux chercheurs de pratiquer une \og lecture distante \fg{} sur de vastes corpus iconographiques.

Dans notre deuxième partie, nous avons commencé par la création d'une ontologie pour entraîner un modèle de vision adapté à des documents historiques. Cette approche, partiellement inspirée de SegmOnto et du dataset S-VED, est pensée pour être compatible avec les contraintes techniques des modèles \gls{yolo} tout en englobant manuscrits et imprimés. Ensuite, nous avons présenté la construction du jeu de données GenHisDoc. Pour constituer ce corpus de référence, le travail s'est appuyé sur l'agrégation de jeux de données libres existants (S-VED, IlluHisDoc, Horae) et sur la récupération d'annotations directement extraites des projets de la plateforme AIKON via l'\gls{api} \gls{iiif}. Pour cela, nous avons passé par une phase d'harmonisation technique et d'harmonisation ontologique suivie d'une phase de correction et d'enrichissement en utilisant le logiciel Label Studio.  Enfin, nous avons disséqué le protocole d'entraînement et d'évaluation de notre travail. Nous avons évalué des points forts (la détection d'estampilles) et des éléments à retravailler (la détection d'ornementations) ainsi que toutes les métriques qui y sont liées. 

% Élargissement  

Nous attirons votre attention sur le fait que les choix des modèles à utiliser ont plus été influencé par une considération logistique (ce qui est déjà utilisé sur AIKON) qu'un réel intérêt technique. Les modèles d'une même génération vont plus ou moins tous être correctement utilisables pour travailler sur des documents historiques et ils auront globalement tous les mêmes comportements. Cependant, les récentes avancées dans le domaine des modèles vision-langage ouvrent la voie à des méthodes dont les capacités, les coûts et les inconvénients sont encore difficilement évaluables. 

La vision apporte beaucoup au travail scientifique : elle permet de simplifier l'heuristique historique en traitant des masses documentaires. Cependant, elle demande une connaissance fine à la fois des documents que l'on veut traiter et du fonctionnement interne de la vision artificielle et de ses métriques, car les biais présents dans les données peuvent venir autant de la sélection du corpus, d'erreurs dans les annotations ou d'un mauvais choix lors de l'évaluation. J'espère que ce travail démontre que les ingénieurs ne remplaceront pas les historiens, mais qu'au contraire c'est un travail de mise en commun des savoirs et de définition des besoins qui permet de structurer la démarche au départ et d'obtenir un résultat pertinent à l'arrivée.